\section*{Введение}
\addcontentsline{toc}{section}{Введение}

\textbf{Актуальность темы.} Проблема «живых» очередей на кафедрах технических вузов, 
таких как Московский Политех, остается крайне острой из-за высокой плотности учебного 
графика и сложности принимаемых работ (лабораторные работы, курсовые проекты). Существующие 
способы организации очереди — физическое присутствие в коридоре, Google-таблицы, 
Telegram-боты — либо неудобны, либо не обеспечивают необходимого уровня безопасности и 
аналитики. В связи с этим актуальной является задача разработки специализированного 
веб-приложения, которое позволит автоматизировать процесс управления очередями, учитывая 
специфику учебного процесса и потребности как студентов, так и преподавателей.

\textbf{Степень разработанности проблемы.} В теории массового обслуживания известны классические 
алгоритмы (FIFO, SJF, Round Robin), однако их применение в образовательной среде имеет особенности, 
связанные с необходимостью учёта успеваемости, возможности доработки задания без потери места 
и сбора статистики. Существующие программные решения (Google Sheets, Telegram-боты) не предоставляют 
встроенных механизмов для такого учёта. Таким образом, создание адаптивного алгоритма очереди с возможностью
 динамического приоритизации и аналитики является актуальной исследовательской задачей.

\textbf{Объект исследования} — процесс организации академического времени и взаимодействия
 «преподаватель-студент» во время защиты лабораторных работ.

\textbf{Предмет исследования} — алгоритмы и программные средства автоматизации управления очередями в режиме реального времени.
